\documentclass[11pt]{article}
\usepackage{icmlko}
\begin{document}

\title{옛 반증이 틀린 게 아니라 그 조건이 사라졌다 --- 고치고 나니 뒤집을 것이 없어졌다}
\author{월드모델 실험실}
\date{노트 595 · 2026년 8월 5일}
\maketitle

\begin{abstract}
노트 $587$ 이 짝 라벨 $k$ 건으로 \textbf{무게}를 맞춰 앱을 올렸다. 그런데
노트 $582$ 는 \textbf{부호}로 아무것도 못 올렸다 --- \emph{다른 시점,
다른 프로토콜}이었다. 같은 바탕 $\cdot$ 같은 $k$ $\cdot$ 같은 씨앗으로
붙였다.

\begin{center}
\begin{tabular}{lrrr}
\hline
$k{=}40$ & A 보정 없음 & \textbf{B 무게} & C 부호 \\
\hline
KR 만화 & $+0.6828$ & $+0.6794$ & $+0.6850$ \\
\textbf{비게임 앱} & $+0.5027$ & $\mathbf{+0.5653}$ & $\mathbf{+0.4792}$ \\
CN 만화 & $+0.3237$ & $+0.2771$ & $\mathbf{+0.1231}$ \\
\hline
\end{tabular}
\end{center}

$\mathrm{B}-\mathrm{C}$ 가 앱 $\mathbf{+0.0860}$ $[+0.0470,+0.1250]$ $\cdot$
CN $\mathbf{+0.1540}$ 로 $0$ 밖 양수다. \textbf{무게가 부호를 이긴다.}

그런데 사전 등록 ②가 \textbf{반증됐다} --- 노트 $463$ 은 ``부호 뒤집기가
KR 을 $+0.6115 \to +0.3833$ 로 해친다''를 쟀는데, 여기서는
$\mathrm{C}-\mathrm{A} = \mathbf{+0.0022}$ 다.
\end{abstract}

\begin{figure}[t]
\centering
\includegraphics[width=\textwidth]{figs/f1.pdf}
\caption{\textbf{가} --- 무게가 부호를 이긴다.
\textbf{나} --- 많이 뒤집을수록 크게 다친다.
\textbf{다} --- 옛 반증의 조건이 사라졌다.}
\end{figure}

\section{왜 463 이 재현되지 않나}

표에 답이 있다. C 가 실제로 \textbf{뒤집은 축 수}를 같이 찍었다.

\begin{center}
\begin{tabular}{lrr}
\hline
$k{=}40$ & 뒤집은 축 수 & $\mathrm{C}-\mathrm{A}$ \\
\hline
KR 만화 & $\mathbf{0.3}$ 개 & $+0.0022$ \\
비게임 앱 & $0.7$ 개 & $-0.0235$ \\
CN 만화 & $\mathbf{1.2}$ 개 & $\mathbf{-0.2006}$ \\
\hline
\end{tabular}
\end{center}

\begin{quote}
\textbf{노트 $586$ 채택으로 방향이 서서 뒤집을 것이 없어졌다.} 노트 $463$ 은
\emph{채택 전} 세대였고, 그때는 하네스 열이 다섯 도메인에서 뒤집혀 있어
``학습 부호와 반대''라는 규칙이 자주 발동했다.
\end{quote}

\textbf{$463$ 이 틀린 게 아니다.} 그 노트가 잰 것은 참이었고, \textbf{그것을
참이게 만든 조건을 우리가 없앴다.}

\section{그래도 부호는 켜지 않는다}

$\mathrm{C}-\mathrm{A}$ 가 앱 $-0.0235$, \textbf{CN $-0.2006$} 이다. 그리고
많이 뒤집는 짝일수록 크게 다친다.

\begin{quote}
CN 은 \textbf{축이 셋뿐}이라 하나를 잘못 뒤집으면 입력의 \textbf{$3$ 분의
$1$} 이 뒤집힌다. $k{=}40$ 에서 평균 $1.2$ 개를 뒤집고 $-0.2006$ 을 잃는다.
\end{quote}

\textbf{켤 근거는 여전히 없다} --- 근거가 바뀌었을 뿐이다. 예전에는
``KR 을 크게 해쳐서''였고 지금은 ``도움이 안 되면서 앱 $\cdot$ CN 을
해쳐서''다.

\section*{무엇이 다른가 --- 손대는 자리}

미리 적어 둔 한계가 그대로 산다. \textbf{B 는 출력을 뒤에서 합치고,
C 는 입력을 앞에서 뒤집는다.}

\begin{quote}
그래서 정확히는 ``무게가 부호보다 낫다''가 아니라 \textbf{``이 두 구현 중
B 가 낫다''}이다. 갈라내려면 셋째 팔(입력을 \emph{무게로} 다시 스케일)이
필요한데 안 만들었다. \textbf{안 만든 것을 안 만들었다고 적는다.}
\end{quote}

\section*{조항}

\begin{quote}
\textbf{옛 반증을 인용할 때는 그 조건이 아직 있는지 본다.} 이 판의 노트는
$600$ 개에 가깝고, 그중 상당수가 \emph{그때의 판} 위에서 참이었다.
\textbf{채택은 갈래를 열기도 닫기도 한다}(노트 $524$)에 하나를 보탠다 ---
\textbf{채택은 옛 반증을 무효로 만들기도 한다.} 그때 할 일은 그 반증을
지우는 것이 아니라 \textbf{조건이 사라졌다고 적는 것}이다.
\end{quote}

이번에는 그것을 \emph{자동으로} 알았다 --- C 가 뒤집은 축 수를 같이
찍었기 때문이다. \textbf{손잡이가 몇 번 발동했는지를 결과와 같이 적으면,
그 손잡이가 왜 안 듣는지가 표에서 보인다.}

\section*{남은 것}

\begin{itemize}
\item \textbf{셋째 팔} --- 입력을 무게로 다시 스케일하는 구현. ``출력 대
입력''과 ``무게 대 부호''를 갈라낸다.
\item \textbf{CN 은 무게도 손해다}($k{=}40$ $-0.0466$). $352$ 행이라 $k$ 를
빼면 채점 집합이 크게 줄고 축이 셋뿐이라 능형이 쓸 것이 적다.
\textbf{작은 짝에는 $k$ 를 쓸 수 없다}는 뜻인데, 하필 작은 짝이 $k$ 를
제일 필요로 한다.
\item \textbf{\texttt{MIN\_K} 를 짝 크기로} --- 지금은 $40$ 고정이다.
CN($352$ 행)에서는 $40$ 도 $11\%$ 다.
\end{itemize}

\end{document}
